
\section{The Setting and Data}
\label{section:data}

Ethiopia installed ethnic federalism in the early 1990's with a coalition of ethnic liberation movements forming a government, changing the unitary nature of the state \citep{Henze1992}. The country is now a federation of different ethnic groups, and the three main groups - Amhara, Oromo, and Tegaru – together constitute about 70percent of the country’s population.\footnote{ estimates are based on the Ethiopian Census 2007.} 

Before the change, Amharic was the dominant language in the country, serving as the official working language of the state and the language of instruction in all primary schools \citep{Heugh2007}. In high schools, English was used as a medium of instruction for all subjects and Amharic was taught as a separate subject. The learning and dissemination of Amharic was enforced by state policy and proficiency in the language was essential for success in the job market.

The new government issued a new language of instruction policy in 1994, empowering the federated ethnic states to establish mother tongue instruction in primary and middle schools under their jurisdiction \citep{Heugh2007}. The states adopted the new language of instruction policy differently, with some implementing it fully (up to and including the first eight grades of schooling, consistent with the policy), and others only partially, switching to English language instruction before students reach high school. A review of the MTI adoption processes in each state, as detailed in \citet{Heugh2007}, shows that the geographical spread of MTI was set randomly and politicians made decisions arbitrarily based on unsubstantiated claims.

Decisions not to fully comply with the national MTI policy in the state of Amhara and the Regional State of Southern Nations and Nationalities (the South) were pushed by local politicians and elites, against the recommendations of local experts in education \citep{Heugh2007}. Indeed, policy makers in these states did not present compelling arguments, justifying how instituting English as a language of instruction in elementary schools would facilitate learning in communities where English is barely spoken and understood, and educators experience well-documented difficulties in communicating content in English. According to \citep{Heugh2007},

\begin{quote}
    ...[M]any educators [in the Amhara state] defended the use of mother tongue to the end of primary schooling...the elite, however, were anxious for English to begin earlier because it becomes the mother of instruction (MOI) in grade nine. Reference was made to Addis Ababa, where English [had] recently been made [medium of instruction] for grades 7 and 8, and (mistakenly) to Tigray region which [had] not in fact adopted English at this level. None of those interviewed were aware of any studies done on what parents want in terms of the mother tongue or how the change might impact on student performance at upper primary levels \citep[][p.~63]{Heugh2007}.
\end{quote}

%Likewise, it is difficult to argue that the choice made in Oromia and Tigray to fully comply with the national policy was motivated exclusively by the objective of promoting the educational outcomes of students. It is plausible that policy makers in these states also had as their objective the development of Afaan Oromo and Tigriyna, languages that  played minimal roles in the country's social, economic and political affairs before the introduction of ethnic federalism in the country. Given the history of contest for political power between Ethiopia's three major ethnic groups -- the Amhara, the Oromo and the Tegaru -- it would be safe to assume that a fuller implementation of the MTI policy in Oromia and Tigray might not have been motivated solely by the goals of enhancing learning.


%\subsection*{Data and Descriptive Figures}

Our data comes from the Ethiopian Young Lives Survey (YLS)\footnote{The YLS is a longitudinal survey of 12,000 children from four developing countries: Ethiopia, India, Peru and Vietnam, curated by the Department of International Development at the University of Oxford.}, and includes a random sample of about 800 youths (the Older Cohort) and their families – in four different regional states of Ethiopia (Amhara, Oromia, The South, and Tigray) as well as the capital, Addis Ababa. The sampled households were surveyed in five rounds, beginning in 2001 when the subjects were eight-years old, with the last round of the survey conducted in 2016 when the subjects were in some form of employment or seeking employment. The Survey includes data on test scores, wage and salary employment, as well as information pertaining to the characteristics of students and their families.

Table \ref{tab:sum-stat} presents descriptive statistics of the variables used in the regression analysis. The first three columns present summary statistics of the sample from the four regions, and the last three columns describe the sample from Addis Ababa.  The table shows that Addis Ababa students have higher test scores compared to the rest of the country. In addition, the employment outcomes are better in Addis Ababa, which is in line with our expectations. Focusing on the sample from the four regional states, we observe that about half had some kind of wage employment, whereas around 16\% were engaged in salaried employment by the fifth round. According to the overall distribution of the Ethiopian population, approximately 73\% of the overall sample live in rural areas.

Figure \ref{fig:employ_gr} displays the percentage of pupils who were employed by the last round of the survey. The region of Tigray outperforms the other regions in terms of both employment outcomes, closely followed by Oromia. The South and the Amhara state appear to have the lowest wage and salary employment rates, with the Amhara region and the South performing the worst in terms of salary employment and wage employment, respectively. It is notable that the two regional states with higher rates of salary and wage employment, Tigray and Oromia, are those that instituted the 1994 MTI policy fully, whereas the underperformers, Amhara and the South, implemented MTI in their curriculum only partially. 
%Figure 2, on the other hand, is a kernel density estimate of standardized maths and language test scores in each region. Standardized language test score displays more variability across the regions, with students in Tigray and Amhara performing better than their counterparts in SNNP and Oromia. In contrast, the standardized maths test scores appear to be more stable.



